Central chiller plants remain major electricity consumers in continuously operated facilities, yet much of the recent literature emphasizes simulator-specific tuning or model predictive control rather than a transparent planning-oriented optimization scaffold. This paper develops a reproducible day-ahead Pyomo mixed-integer linear programming framework for a heterogeneous multi-chiller plant with chilled-water storage, chilled-water pumps, condenser-water pumps, and cooling-tower auxiliaries. Beyond unit commitment, piecewise-linear efficiency curves, and storage dynamics, the formulation explicitly robustifies three uncertainty channels: cooling load, electricity price, and wet-bulb-driven efficiency/capacity degradation. The last component is introduced through conservative but linear ambient-aware chiller-power coefficients and wet-bulb-dependent available-capacity limits, so the resulting model remains MILP-solvable. We evaluate the framework on a semi-synthetic 24-hour benchmark with three large chillers, one small chiller, time-of-use pricing, and weather-linked stress tests. Relative to the lowest-cost deterministic schedule with storage, the proposed ambient-aware robust schedule increases nominal operating cost from 5584.7 to 6035.6 while reducing the out-of-sample mean cooling shortfall from 727.4~kWh to 61.7~kWh and the correlated-stress mean shortfall from 651.1~kWh to 35.4~kWh. Relative to a load-and-price-robust schedule with storage, the ambient-aware extension further lowers correlated-stress shortage from 36.8 to 35.4~kWh and correlated out-of-sample energy cost from 6052.6 to 6049.6. We therefore position the work as a reproducible full-plant benchmark for ambient-aware robust chiller scheduling rather than a field-validated deployment study.
